\section{Tugas Scanner dan Antarmuka Lexer}
\label{sec:ch05-peran}

\subsection{Posisi Lexer dalam Pipeline}

Lexer, atau \eng{scanner}, adalah fase pertama kompilator. Ia membaca
aliran karakter dari program sumber dan mengelompokkannya menjadi
token. Token adalah unit bermakna terkecil yang dikenali parser.
Lexer membuang spasi, baris baru, dan komentar, sehingga parser tidak
perlu menangani detail leksikal.

Guntoro merumuskan tugas scanner sebagai rangkaian kerja yang jelas:
membaca program sumber karakter demi karakter, membuang spasi dan
komentar, mengelompokkan karakter menjadi token, menyerahkan token
ke parser, dan mencatat identifier ke tabel simbol bila diperlukan.
Pembahasan berikut mengikuti urutan tugas tersebut.

\begin{figure}[H]
\centering
\begin{tikzpicture}[node distance=6mm and 9mm]
  \node[data, text width=2.2cm] (src) {Program\\sumber};
  \node[phase, text width=2.4cm, right=of src] (lex) {Lexer};
  \node[data, text width=2.4cm, right=of lex] (tok) {Aliran\\token};
  \node[phase, text width=2.4cm, right=of tok] (par) {Parser};
  \node[phase, text width=2.2cm, below=of lex] (sym) {Tabel\\simbol};
  \draw[arrow] (src) -- (lex);
  \draw[arrow] (lex) -- (tok);
  \draw[arrow] (tok) -- (par);
  \draw[arrow, dashed] (lex) -- (sym);
  \draw[arrow, dashed] (par) -- (sym);
\end{tikzpicture}
\caption{Antarmuka lexer, token stream, parser, dan tabel simbol}
\label{fig:ch05-antarmuka}
\end{figure}

\subsection{Antarmuka dengan Parser}

Antarmuka yang paling umum adalah \eng{pull model}: parser memanggil
fungsi \code{next\_token()} setiap kali membutuhkan token berikutnya.
Lexer mempertahankan posisi baca dan mengembalikan token satu per
satu. Beberapa lexer menyediakan \code{peek()} untuk melihat token
berikutnya tanpa mengonsumsinya.

Pada \eng{push model}, lexer aktif memanggil parser atau mengisi
buffer token. Model ini jarang dipakai pada kompilator sederhana,
tetapi muncul pada sistem \eng{streaming} atau \eng{event-based}.
Untuk proyek Subset-C, kita memakai pull model karena lebih mudah
diuji dan di-debug.

\subsection{Kontrak Lexer}

Lexer memiliki kontrak berikut:
\begin{itemize}[leftmargin=*]
  \item mengembalikan token berikutnya, atau \code{TOK\_EOF} pada
        akhir masukan;
  \item mengisi atribut token: jenis, lexeme, baris, kolom, dan nilai
        bila ada;
  \item melewati spasi, tab, baris baru, dan komentar;
  \item menghasilkan \code{TOK\_ERROR} untuk karakter ilegal tanpa
        menghentikan seluruh proses;
  \item tidak mengubah program sumber.
\end{itemize}

Kontrak ini penting karena parser dan pengujian bergantung pada
perilaku yang konsisten. Jika lexer diam-diam menghentikan proses
pada kesalahan pertama, parser tidak dapat melakukan pemulihan.

\begin{infobox}
Lexer tidak perlu memahami makna program. Tugasnya hanya mengenali
bentuk token. Keputusan seperti ``apakah variabel sudah
dideklarasikan?'' adalah tugas analisis semantik, bukan lexer.
\end{infobox}
